\section{A first attempt: kernel reweighting}

An initial idea to recover submodularity was to incorporate the weights directly into the kernel by defining
\begin{equation*}
k_w(x,x')
=
w(x)\,k(x,x')\,w(x').
\end{equation*}
The corresponding information gain then becomes
\begin{equation*}
\Gamma_T(k_w)
=
\max_{x_1,\dots,x_T} \frac{1}{2}
\log
\det
\left(
I+\sigma^{-2}K_T^{(w)}
\right),
\end{equation*}
where
\begin{equation*}
K_T^{(w)}(i,j)
=
w(x_i)k(x_i,x_j)w(x_j).
\end{equation*}

This restores submodularity, since the objective is again ordinary Gaussian process information gain. However, the Gaussian process prior itself has changed. The posterior variance now satisfies
\begin{equation*}
\sigma_{n,w}^2(x)
=
k_w(x,x)
-
k_w(x,X_n)
\left(
K_n^{(w)}+\sigma^2 I
\right)^{-1}
k_w(X_n,x),
\end{equation*}
which is not equal to
\begin{equation*}
w(x)^2\sigma_n^2(x).
\end{equation*}
Thus the weighting modifies the geometry of the Gaussian process itself rather than only the acquisition objective. This leads to the assumption that higher-weighted regions have inherently larger variance (amplitude of $f$), changing the posterior itself and altering predicitions. Therefore, the resulting optimisation problem is fundamentally different from the original weighted formulation and we are no longer solving the original problem. 